% 来源：成文/A\_05\_问题分析.md
\section{问题分析}
中药材的干燥是决定药材品质的关键工序，热风烘干是产地常用的方式。就数学形式而言，题目要求求解的是一类以药材内部温度场与干基含水率场为状态变量的非稳态传热传质耦合问题，圆柱几何、时变环境边界与随含水率变化的物性共同决定了它的耦合结构。

就方程类型而言，这是抛物型偏微分方程组：时间上只含一阶导数、空间含二阶导数，故定解条件需且仅需一套初值与边界条件。就任务性质而言，四问分两类：问题一、问题二为正演，问题三、问题四为反演，由达标判据反求达标时刻；四问均无目标函数与决策变量，故不宜套用优化模型。全文难点在物性使两场牵制、低含水率端经验式大幅外推、界面系数取法决定反演答案三处。

四问共用同一个物理过程：圆柱形药材在热风环境中，内部以导热方式升温、以扩散方式失水，两者通过物性相互影响。因此全文只需建立一个模型，四问的区别仅在于四组旋钮的位置：物性体系、耦合方式、时程与判据、几何是否收缩。本文采用的技术路线是从守恒律建立控制方程，以元体平衡（单元中心有限体积）离散，以分场选取的时间格式推进，再用解析解对拍、独立实现互验与守恒核算三类手段验证，并以灵敏度分析界定结论的稳健范围。

本文输入全部来自题面与附件。环境温度与水分浓度取自附件 1（241 点、间隔 60 s），药材半径取自附件 2（145 点、间隔 1800 s），均无缺失与离群点，节点间线性插值、不作平滑；附件不覆盖全程，恒温段按末段实测均值延拓、自问题二起生效，物性与表面系数、达标阈值均由题面给定。

\subsection{问题一的分析}
问题的物理场景是预热平衡阶段：环境温度尚在上升，水分尚未大量脱除。关键判断是能否用集总参数法。传热毕奥数 $Bi_T=hR_0/k$ 约为 1.39、传质毕奥数 $Bi_m=k_mR_0/D$ 约为 3.24，二者均为一阶量级，说明内部导热阻与外部对流阻同量级，药材内部必然存在显著径向梯度，集总参数法不适用，必须做分布式求解\textsuperscript{\cite{ref1,ref8}}。按判据分档可写得更清晰：毕奥数小于 0.1 时内部阻力可略、可用集总参数法，介于 0.1 与 100 之间时须作分布式求解，大于 100 时控制权转入内部；本例两数均落中间档\textsuperscript{\cite{ref1,ref8}}。随干燥推进与半径收缩，传质毕奥数下降，控制权由外部对流向内部扩散迁移。

第二个判断是能否用半无限介质近似：由于扩散域径向厚度与半径之比不可忽略，忽略曲率的半无限近似会低估径向扩散路径。本文以对照计算验证：同一条件下半无限近似把表面含水率高估约 6.7\%（$1.6113$ 对 $1.5104$ kg/kg），故必须保留圆柱坐标下的曲率项。

第三个判断是两场能否解耦：问题一的物性为常值，扩散系数只依赖含水率而不依赖温度，含水率方程也不含温度，两场之间没有反向依赖，故可严格解耦、依次求解——这是本文在问题一能对温度场使用时间方向精确推进的前提。

\subsection{问题二的分析}
与问题一相比，变化集中在两处：物性从常值升为含水率与温度的函数（附录 3），两场由解耦转为双向强耦合\textsuperscript{\cite{ref2,ref9}}。因此必须回答为何必须联立迭代：若先算温度再算含水率而不回代，物性始终取旧状态，而干燥中段物性变化最快，误差会沿时间累积。本文采用交替推进加物性逐内部子步更新的方式化解，子步取 $1/32$ s 时物性冻结造成的滞后不进入四位小数量级。

\subsection{问题三的分析}
与问题二相比，仅动两处：时程由固定 3 h 改为事件驱动（达到判据即停），并新增终止判据（全域含水率严格小于 $0.15$ kg/kg）。分析要点有二：其一，中心恒为最慢点（径向扩散路径最长、边界通量最先影响表面），故"全域达标"与"中心达标"等价，判据可只在中心监测；其二，几何半径在本题设定下全程固定（收缩是问题四的设定），故模型与离散可完全复用问题二。这决定了本文对问题三的处理是"复用加终止机制"，而不是另建模型。

\subsection{问题四的分析}
在问题三基础上叠加尺寸收缩（移动边界）与附录 4 物性。分析要点是为何必须做坐标变换：若在固定坐标系中求解，计算域随半径收缩而变小，网格点与输出位置会逐渐错位，表面附近的控制体不断被压缩，既引入插值误差又恶化表面列的收敛性。本文引入物料坐标 $\xi=r/R(t)$，把收缩域映射为固定计算域 $[0,1]$，使表面恒为节点、输出点恒为计算点，从源头消除上述两类困难。还需说明该变换的前提：附件 2 只给出半径的单一序列，药材内部各处的收缩分配未被规定，故本文按仿射收缩处理，并在第三章假设中一并声明其依据与适用前提。

四问的要素对照见下表。



\settabw{10}
\begin{longtable}{>{\raggedright\arraybackslash}m{\dimexpr 0.138\tblw\relax}>{\centering\arraybackslash}m{\dimexpr 0.209\tblw\relax}>{\centering\arraybackslash}m{\dimexpr 0.213\tblw\relax}>{\centering\arraybackslash}m{\dimexpr 0.209\tblw\relax}>{\centering\arraybackslash}m{\dimexpr 0.231\tblw\relax}}
  \caption{四问要素对照（同一模型 × 四组旋钮）}\label{tab:2-1}\\
  \toprule
    要素 & 问题一 & 问题二 & 问题三 & 问题四 \\
  \midrule
  \endfirsthead
  \multicolumn{5}{r}{\footnotesize（续）}\\
  \toprule
    要素 & 问题一 & 问题二 & 问题三 & 问题四 \\
  \midrule
  \endhead
  \bottomrule
  \endfoot
  \bottomrule
  \endlastfoot
  物性体系 & 附录 2（常物性） & 附录 3（变物性） & 附录 3（变物性） & 附录 4（变物性） \\
  耦合方式 & 严格解耦（依次求解） & 双向强耦合（交替推进） & 双向强耦合 & 双向强耦合 \\
  时程与终止 & 固定 $0$–$1800$ s & 固定 $0$–$3$ h & 事件驱动（达标即停） & 事件驱动（达标即停） \\
  几何 & 固定半径 $R_0$ & 固定半径 $R_0$ & 固定半径 $R_0$ & 收缩 $R(t)$，物料坐标 \\
  待求量 & 两场分布 & 两场分布 & 达标时间 & 烘干时长 \\
\end{longtable}

四问由此形成"框架 → 耦合 → 反演 → 变域"的递进闭环：问题一建立解耦情形下的完整求解与验证框架；问题二在同一方程形式下引入变物性与双向耦合，并给出与问题一重叠段差异的归因；问题三复用问题二的求解核，加入反演机制得到达标时间；问题四再叠加移动边界，并用问题三的固定半径结果作为对照臂分解收缩效应。四问的旋钮分配与递进关系，也决定了后文各章的呈现次序与检验重点。
